\documentclass[11pt]{article}

\usepackage[margin=1.05in]{geometry}
\usepackage[T1]{fontenc}
\usepackage[utf8]{inputenc}
\usepackage{lmodern}
\usepackage{amsmath,amssymb,amsthm,mathtools}
\usepackage{microtype}
\usepackage[colorlinks=true,linkcolor=blue,citecolor=blue,urlcolor=blue]{hyperref}

\newtheorem{theorem}{Theorem}
\newtheorem{corollary}{Corollary}
\newtheorem{lemma}{Lemma}
\newtheorem{problem}{Problem}
\newtheorem{remark}{Remark}

\DeclareMathOperator{\Ree}{Re}
\newcommand{\C}{\mathbb C}

\title{A Negative Answer to the Universal-Function Version of Erd\H{o}s's Third Question in Problem 514}
\author{}
\date{Revised draft, May 2026}

\begin{document}
\maketitle

\begin{abstract}
Let \(f\) be a transcendental entire function and let
\[
M_f(r)=\max_{|z|=r}|f(z)|.
\]
Erd\H{o}s asked, among other things, whether there exists a path to infinity
along which \(|f(z)|\) tends to infinity faster than a fixed function of
\(M_f(|z|)\), for instance faster than a positive power of \(M_f(|z|)\).
This note addresses only that third question, in the natural universal sense
that the comparison function is fixed in advance and is independent of \(f\).
We prove a negative answer in a strong form.  For every nondecreasing function
\(\Phi(T)\to\infty\) there is a transcendental entire function \(f\) such that
every path to infinity \(\gamma\) satisfies
\[
\liminf_{t\to\infty}
\frac{|f(\gamma(t))|}{\Phi(M_f(|\gamma(t)|))}=0.
\]
Consequently no universal comparison function, even one growing much more
slowly than every positive power, can be forced in the universal-comparison
interpretation of Erd\H{o}s's third question.
The result follows from a theorem of Hayman on the growth of entire functions
on asymptotic paths, after a short selection argument for Hayman's auxiliary
growth function.  We also include a direct construction, which gives a compact
and formalization-friendly alternative proof of the same theorem.
\end{abstract}

\section{The formulation addressed here}

Let \(f\) be a transcendental entire function, that is, an entire function
which is not a polynomial.  For \(r>0\), put
\[
M_f(r):=\max_{|z|=r}|f(z)|.
\]
A path to infinity is a continuous map
\[
\gamma:[0,\infty)\to\C
\]
such that
\[
|\gamma(t)|\to\infty \qquad (t\to\infty).
\]

Erd\H{o}s's third question in Problem 514 asks whether one can force
\(|f(z)|\) to tend to infinity along a path faster than a fixed function of
\(M_f(r)\); see \cite[p. 249]{Erdos1961} and the later problem list
\cite{Erdos1982}.  We formalize the relevant positive assertion as follows.

\begin{problem}[Universal comparison-function formulation]
Does there exist a function \(\Psi:[T_0,\infty)\to(0,\infty)\), with
\(\Psi(T)\to\infty\) as \(T\to\infty\), independent of \(f\), such that every
transcendental entire function \(f\) has a path to infinity \(\gamma\) for which
\[
\frac{|f(\gamma(t))|}{\Psi(M_f(|\gamma(t)|))}\to\infty
\qquad (t\to\infty)?
\]
\end{problem}

All statements involving \(\Psi(M_f(r))\) or \(\Phi(M_f(r))\) are understood for
sufficiently large \(r\), so that \(M_f(r)\ge T_0\).  This is harmless, since
\(M_f(r)\to\infty\) for every nonconstant entire function.

It is no loss to restrict attention to nondecreasing comparison functions.
Indeed, suppose that a not-necessarily-monotone function
\(\Psi:[T_0,\infty)\to(0,\infty)\), with \(\Psi(T)\to\infty\), were to work.
Choose \(T_1\ge T_0\) so large that \(\Psi(S)\ge 1\) for all \(S\ge T_1\), and
on \([T_1,\infty)\) define the lower envelope
\[
\Phi(T):=\inf_{S\ge T}\Psi(S).
\]
Then \(1\le\Phi\le\Psi\), the function \(\Phi\) is nondecreasing, and
\(\Phi(T)\to\infty\).  The last assertion follows because, for each \(C>0\),
we have \(\Psi(S)\ge C\) for all sufficiently large \(S\).  A path with
\(|f(\gamma(t))|/\Psi(M_f(|\gamma(t)|))\to\infty\) would then also have
\(|f(\gamma(t))|/\Phi(M_f(|\gamma(t)|))\to\infty\).
Thus it suffices to rule out nondecreasing \(\Phi\)'s; discarding the initial
interval \([T_0,T_1)\) is irrelevant for all asymptotic statements.

\paragraph{Related work.}
The first two questions in Problem 514 are closely related to the
subharmonic-function path results of Huber, Lewis--Rossi--Weitsman, and Wu.
Chojecki's note \cite{Chojecki2026} applies this line to obtain the existence
and length assertions.  For the parenthetical power-type part of the third
question, Chojecki used Langley's formulation \cite[Theorem~1.4]{Langley2022}
of a consequence of Barth--Brannan--Hayman \cite{BarthBrannanHayman1978}.
There is also a complementary finite-lower-order/multiple-tract line: Rossi and
Weitsman \cite{RossiWeitsman1985} proved a power lower bound in \(|z|\) along
an asymptotic path when a high-level set has at least two components, and Toda
\cite{Toda1993} sharpened this type of result, including a subharmonic
analogue.  These results are positive path-growth theorems in terms of \(|z|\)
and tract geometry, rather than universal lower bounds in terms of \(M_f(r)\).
The present note is concerned only with the universal-comparison formulation;
for that formulation Hayman's theorem gives the result after Lemma
\ref{lem:hayman-selection}, and the direct construction below is a
self-contained and formalization-friendly proof of the same corollary.

We prove the following stronger negative statement.

\begin{theorem}\label{thm:main}
Let
\[
\Phi:[T_0,\infty)\to(0,\infty)
\]
be nondecreasing and satisfy \(\Phi(T)\to\infty\) as \(T\to\infty\).  Then there
exists a transcendental entire function \(f\) such that for every path to
infinity \(\gamma\),
\[
\liminf_{t\to\infty}
\frac{|f(\gamma(t))|}{\Phi(M_f(|\gamma(t)|))}=0.
\]
Since the quotient is nonnegative whenever it is defined, the conclusion is
equivalently the following epsilon form: for every path \(\gamma\), every
\(\eta>0\), and every \(T_\ast\ge 0\), there exists \(t\ge T_\ast\), with
\(M_f(|\gamma(t)|)\ge T_0\), such that
\[
\frac{|f(\gamma(t))|}{\Phi(M_f(|\gamma(t)|))}\le \eta.
\]
\end{theorem}

\begin{corollary}\label{cor:third-question}
There is no universal function \(\Psi(T)\to\infty\), independent of \(f\), for
which every transcendental entire function has a path to infinity along which
\[
|f(z)|/\Psi(M_f(|z|))\to\infty.
\]
In particular, the parenthetical power-type case \(\Psi(T)=T^\varepsilon\),
\(\varepsilon>0\), also fails.
\end{corollary}

\section{Hayman's theorem and the selection step}

The main input is Hayman's theorem on the growth of entire functions on
asymptotic paths.  We use ``integral function'' in Hayman's classical sense,
namely entire function.  A path \(\Gamma\) is an asymptotic path of \(f\) if
\(z\to\infty\) along \(\Gamma\) and \(f(z)\to\infty\) along \(\Gamma\).

\begin{theorem}[Hayman {\cite[Theorem 2, p.~253]{Hayman1960}}]\label{thm:hayman}
Let \(\lambda(r)\) be a positive increasing function of \(r>0\) such that
\[
\frac{\log \lambda(r)}{\log r}\to\infty
\qquad (r\to\infty).
\]
Then there exist two entire functions \(f_1\) and \(f_2\), each of infinite
lower order, such that, for all sufficiently large \(r\),
\[
\log M(r,f_1)<\exp(\lambda(r))
\qquad\text{and}\qquad
\log M(r,f_2)>\exp(\lambda(r)).
\]
Moreover, for \(j=1,2\), if \(\Gamma_j\) is an asymptotic path of \(f_j\), then
\[
\limsup_{\substack{z\to\infty\\ z\in\Gamma_j}}
\frac{\log\log |f_j(z)|}{\log |z|}<\infty .
\]
\end{theorem}

Only the fast-growth function \(f_2\) is needed below.

\begin{lemma}[Selection of Hayman's auxiliary growth]\label{lem:hayman-selection}
Let \(\Phi:[T_0,\infty)\to(0,\infty)\) be nondecreasing and satisfy
\(\Phi(T)\to\infty\) as \(T\to\infty\).  There exists a positive increasing
function \(\lambda(r)\), \(r>0\), such that
\[
\frac{\log \lambda(r)}{\log r}\to\infty
\]
and, for every positive integer \(k\),
\[
\Phi(\exp(\exp(\lambda(r))))\ge \exp(r^k)
\]
for all sufficiently large \(r\).
\end{lemma}

\begin{proof}
For each integer \(n\ge 2\), choose \(S_n\ge \max\{T_0,e\}\) such that
\[
\Phi(S_n)\ge \exp(n^n),
\]
which is possible because \(\Phi(T)\to\infty\).  Choose a strictly increasing
sequence \(a_n\), \(n\ge 2\), such that
\[
a_n\ge n^n,
\qquad
a_n\ge \log\log S_n,
\qquad
\exp(\exp(a_n))\ge S_n.
\]
Define \(\lambda(n)=a_n\) for integers \(n\ge 2\), interpolate linearly on each
interval \([n,n+1]\), and extend \(\lambda\) on \((0,2]\) so that it is positive
and increasing on \((0,\infty)\).

If \(n\le r\le n+1\), then \(\lambda(r)\ge a_n\), and hence
\[
\frac{\log \lambda(r)}{\log r}
\ge
\frac{\log a_n}{\log(n+1)}
\ge
\frac{\log(n^n)}{\log(n+1)}
\to\infty .
\]
Thus \(\lambda\) has Hayman's required growth.

Fix a positive integer \(k\).  For \(n\le r\le n+1\), monotonicity of \(\Phi\)
and the choices above give
\[
\Phi(\exp(\exp(\lambda(r))))
\ge
\Phi(\exp(\exp(a_n)))
\ge
\Phi(S_n)
\ge
\exp(n^n).
\]
For all sufficiently large \(n\), \(n^n\ge r^k\) throughout
\([n,n+1]\).  Therefore
\[
\Phi(\exp(\exp(\lambda(r))))\ge \exp(r^k)
\]
for all sufficiently large \(r\), as required.
\end{proof}

\begin{proof}[Proof of Theorem \ref{thm:main} from Hayman's theorem]
Choose \(\lambda\) as in Lemma \ref{lem:hayman-selection}, and let \(f=f_2\) be the
fast-growth entire function supplied by Theorem \ref{thm:hayman}.  Since \(f\)
has infinite lower order, it is transcendental entire.  Hayman's maximum
modulus conclusion gives, for all sufficiently large \(r\),
\[
M_f(r)>\exp(\exp(\lambda(r))).
\]

Let \(\gamma:[0,\infty)\to\C\) be any path to infinity.  Recall that, by our
definition, this means that \(\gamma\) is continuous and
\(|\gamma(t)|\to\infty\).  If \(\gamma\) is not an asymptotic path of \(f\),
then \(|f(\gamma(t))|\) does not tend to infinity.  Hence there are a constant
\(B>0\) and a sequence \(t_m\to\infty\) such that
\[
|f(\gamma(t_m))|\le B
\]
for all \(m\).  Since \(f\) is nonconstant entire, \(M_f(r)\to\infty\), and
hence
\[
\Phi(M_f(|\gamma(t_m)|))\to\infty .
\]
Along this sequence the quotient in Theorem \ref{thm:main} tends to zero.

It remains to consider the case where \(\gamma\) is an asymptotic path of
\(f\).  Let \(\Gamma=\gamma([0,\infty))\).  The compactness of
\(\gamma([0,T])\) for each fixed \(T\) implies that estimates holding as
\(z\to\infty\) along \(\Gamma\) hold eventually as \(t\to\infty\) along
\(\gamma\).  Hayman's path estimate therefore gives a finite constant \(C\)
such that, eventually along \(\gamma\), with \(r=|\gamma(t)|\),
\[
\log\log |f(\gamma(t))|\le C\log r.
\]
Equivalently, after increasing \(C\) harmlessly and discarding an initial
segment of the path,
\[
|f(\gamma(t))|\le \exp(r^C).
\]
Choose an integer \(k>C\).  By the fast-growth bound for \(M_f\), monotonicity
of \(\Phi\), and Lemma \ref{lem:hayman-selection},
\[
\Phi(M_f(r))
\ge
\Phi(\exp(\exp(\lambda(r))))
\ge
\exp(r^k)
\]
eventually along the path.  Therefore
\[
0\le
\frac{|f(\gamma(t))|}{\Phi(M_f(|\gamma(t)|))}
\le
\exp(r^C-r^k)\to 0 .
\]
In particular the required liminf is zero.  Since \(\gamma\) was arbitrary,
Theorem \ref{thm:main} follows.
\end{proof}

\section{A direct alternative construction}

The proof above shows that Theorem \ref{thm:main} is a corollary of Hayman's
Theorem \ref{thm:hayman} plus the elementary selection lemma.  We now give a
self-contained construction proving the same theorem directly.  The point of
this second proof is not priority for the theorem statement, but rather a short
construction whose estimates are explicit and well suited to formal checking.

Fix \(\Phi\) as in Theorem \ref{thm:main}.  Put
\[
\varepsilon_n:=(-1)^{n+1},\qquad n=1,2,3,\ldots .
\]
Thus \(\varepsilon_n=1\) for odd \(n\) and \(\varepsilon_n=-1\) for even \(n\).

Choose
\[
x_0>0
\]
so large that \(e^x/2\ge T_0\) whenever \(x\ge x_0\).  For \(x\ge x_0\), define
\[
A(x):=\min\left\{\Phi(e^x/2),\frac{e^x}{4}\right\}.
\]
Since both \(\Phi(e^x/2)\) and \(e^x/4\) tend to infinity, we have
\[
A(x)\to\infty \qquad (x\to\infty).
\]

We shall construct positive real numbers
\[
r_n,\qquad X_n,\qquad \lambda_n,\qquad a_n
\]
with \(r_n\to\infty\), \(X_n\to\infty\), and
\[
a_n=\lambda_n r_n-X_n.
\]
The entire function will be
\[
 f(z):=\sum_{n=1}^{\infty}
 \exp(-a_n+\varepsilon_n\lambda_n z).
\]

The parameters are chosen so that the following off-diagonal estimates hold
for all \(m\ne n\):
\begin{equation}\label{eq:offdiag}
\exp(-a_m+\lambda_m r_n)
\le
2^{-m-n-4}A(X_n).
\end{equation}
We also require
\begin{equation}\label{eq:AXn-large}
A(X_n)\ge 2^{n+4},
\qquad
X_n\ge n\log r_n+n,
\qquad
X_n\ge x_0,
\end{equation}
and
\begin{equation}\label{eq:an-large}
a_n\ge n^2+\lambda_n n.
\end{equation}

\begin{lemma}\label{lem:choice}
The sequences \(r_n,X_n,\lambda_n,a_n\) can be chosen so that
\eqref{eq:offdiag}, \eqref{eq:AXn-large}, and \eqref{eq:an-large} hold.
\end{lemma}

\begin{proof}
We argue by induction.  Suppose that the parameters have already been chosen
for indices \(<n\).  Choose
\[
r_n>\max\{r_{n-1},n\},
\]
with the convention that there is no restriction involving \(r_0\) when
\(n=1\).

The quantities
\[
\exp(-a_m+\lambda_m r_n),\qquad m<n,
\]
are now fixed and finite.  Since \(A(x)\to\infty\), we may choose \(X_n\) so
large that \eqref{eq:AXn-large} holds and, for every \(m<n\),
\[
\exp(-a_m+\lambda_m r_n)
\le
2^{-m-n-4}A(X_n).
\]
This gives the off-diagonal estimates in which already constructed terms are
evaluated at the new radius \(r_n\).

It remains to make the new term small at all previous radii.  For each
\(k<n\), the inequality
\[
\exp(X_n-\lambda_n(r_n-r_k))
\le
2^{-n-k-4}A(X_k)
\]
will hold once \(\lambda_n\) is sufficiently large, because \(r_n-r_k>0\).
Since there are only finitely many such \(k\), one value of \(\lambda_n\) can
make all these inequalities true.  Increasing \(\lambda_n\) further if
necessary, we also ensure
\[
\lambda_n(r_n-n)\ge X_n+n^2,
\]
which is possible because \(r_n>n\).  Now set
\[
a_n:=\lambda_n r_n-X_n.
\]
Then
\[
a_n=\lambda_n r_n-X_n\ge n^2+\lambda_n n,
\]
so \eqref{eq:an-large} holds, and for \(k<n\),
\[
\exp(-a_n+\lambda_n r_k)
=
\exp(X_n-\lambda_n(r_n-r_k))
\le
2^{-n-k-4}A(X_k).
\]
Thus all required estimates involving index \(n\) and previous indices hold.
The induction is complete.
\end{proof}

\section{The function is transcendental entire}

Let
\[
 f(z)=\sum_{n=1}^{\infty}
 \exp(-a_n+\varepsilon_n\lambda_n z)
\]
be the function constructed above.

First we show that the series converges locally uniformly.  Fix \(R>0\) and
suppose \(|z|\le R\).  If \(n>R\), then by \eqref{eq:an-large},
\[
-a_n+\lambda_n |z|
\le
-a_n+\lambda_n R
\le
-a_n+\lambda_n n
\le
-n^2.
\]
Hence
\[
\left|\exp(-a_n+\varepsilon_n\lambda_n z)\right|
\le e^{-n^2}
\]
for all sufficiently large \(n\), uniformly on \(|z|\le R\).  Therefore the
series converges uniformly on compact subsets of \(\C\), and \(f\) is entire.

Next we show that \(f\) is not a polynomial.  At the point
\[
z_n:=\varepsilon_n r_n,
\qquad |z_n|=r_n,
\]
the \(n\)-th term equals
\[
\exp(-a_n+\varepsilon_n\lambda_n z_n)
=
\exp(-a_n+\lambda_n r_n)
=e^{X_n}.
\]
For the other terms, \eqref{eq:offdiag} gives
\[
\sum_{m\ne n}
\left|\exp(-a_m+\varepsilon_m\lambda_m z_n)\right|
\le
\sum_{m\ne n}
\exp(-a_m+\lambda_m r_n)
\le
2^{-n-4}A(X_n).
\]
Since \(A(X_n)\le e^{X_n}/4\), this is at most \(e^{X_n}/16\).  Consequently,
\begin{equation}\label{eq:M-lower}
M_f(r_n)
\ge |f(z_n)|
\ge \frac12 e^{X_n}.
\end{equation}
Together with \eqref{eq:AXn-large}, this gives
\[
M_f(r_n)
\ge
\frac12 e^{X_n}
\ge
\frac12 e^n r_n^n.
\]
If \(f\) were a polynomial of degree \(d\), then \(M_f(r)=O(r^d)\), which is
incompatible with the last estimate along the sequence \(r_n\) as \(n\to\infty\).
Thus \(f\) is transcendental entire.

\section{Smallness on the wrong half-plane}

Fix \(n\), and let \(|z|=r_n\).  We call the half-plane
\[
\varepsilon_n\Ree z>0
\]
the correct half-plane for the \(n\)-th circle.  Thus the correct half-plane is
the right half-plane for odd \(n\) and the left half-plane for even \(n\).

Suppose \(|z|=r_n\) and \(\varepsilon_n\Ree z\le 0\).  Then the main term
satisfies
\[
\left|\exp(-a_n+\varepsilon_n\lambda_n z)\right|
=
\exp(-a_n+\varepsilon_n\lambda_n\Ree z)
\le e^{-a_n}
\le 1.
\]
The other terms are bounded by
\[
\sum_{m\ne n}
\left|\exp(-a_m+\varepsilon_m\lambda_m z)\right|
\le
\sum_{m\ne n}\exp(-a_m+\lambda_m r_n)
\le
2^{-n-4}A(X_n).
\]
By \eqref{eq:AXn-large}, \(A(X_n)\ge 2^{n+4}\), and hence
\[
1\le 2^{-n-4}A(X_n).
\]
Therefore
\[
|f(z)|
\le
2^{-n-3}A(X_n).
\]
On the other hand, from \eqref{eq:M-lower} and the monotonicity of \(\Phi\),
\[
\Phi(M_f(r_n))
\ge
\Phi(e^{X_n}/2)
\ge
A(X_n),
\]
where \(e^{X_n}/2\ge T_0\) follows from \(X_n\ge x_0\).  We have proved the key
estimate
\begin{equation}\label{eq:wrong-halfplane}
|z|=r_n,
\quad
\varepsilon_n\Ree z\le 0
\quad\Longrightarrow\quad
|f(z)|
\le
2^{-n-3}\Phi(M_f(r_n)).
\end{equation}

\section{Smallness on the imaginary axis}

For any real \(y\),
\[
\left|\exp(-a_n+\varepsilon_n\lambda_n i y)\right|=e^{-a_n}.
\]
Thus, by \eqref{eq:an-large},
\[
|f(iy)|
\le
\sum_{n=1}^{\infty} e^{-a_n}
\le
\sum_{n=1}^{\infty} e^{-n^2}
=:
B<\infty.
\]
Since \(f\) is transcendental entire, in particular nonconstant, \(M_f(r)\to
\infty\).  Hence \(\Phi(M_f(r))\to\infty\).  It follows that
\begin{equation}\label{eq:imag-axis}
\frac{|f(iy)|}{\Phi(M_f(|y|))}
\to 0
\qquad (|y|\to\infty),
\end{equation}
where the quotient is understood for sufficiently large \(|y|\), so that
\(M_f(|y|)\ge T_0\).

\section{Every path has small values infinitely often}

Let \(\gamma:[0,\infty)\to\C\) be any path to infinity.  Since \(r_n\to\infty\),
for every sufficiently large \(n\) the path meets the circle \(|z|=r_n\).  Let
\[
\tau_n:=\inf\{t\ge 0: |\gamma(t)|=r_n\}
\]
be the first hitting time.  For all large \(n\), \(\tau_n\) is finite,
\(\tau_n<\tau_{n+1}\), and \(\tau_n\to\infty\).  These facts follow from the
continuity of \(|\gamma(t)|\) and the condition \(|\gamma(t)|\to\infty\).

There are two cases.

First suppose that for infinitely many \(n\),
\[
\varepsilon_n\Ree \gamma(\tau_n)\le 0.
\]
For those \(n\), estimate \eqref{eq:wrong-halfplane} gives
\[
\frac{|f(\gamma(\tau_n))|}
{\Phi(M_f(|\gamma(\tau_n)|))}
\le
2^{-n-3}\to 0.
\]
Thus the desired liminf is zero.

It remains to consider the opposite case.  Then, for all sufficiently large
\(n\),
\[
\varepsilon_n\Ree \gamma(\tau_n)>0.
\]
Since \(\varepsilon_n\) changes sign with \(n\), the real parts
\(\Ree \gamma(\tau_n)\) and \(\Ree \gamma(\tau_{n+1})\) have opposite signs for
all large \(n\).  By continuity of \(\gamma\), there exists
\[
s_n\in[\tau_n,\tau_{n+1}]
\]
such that
\[
\Ree \gamma(s_n)=0.
\]
Write \(\gamma(s_n)=iy_n\), where \(y_n\in\mathbb R\).  Since
\(s_n\ge\tau_n\to\infty\) and \(|\gamma(t)|\to\infty\), we have
\(|y_n|=|\gamma(s_n)|\to\infty\).  Hence \eqref{eq:imag-axis} gives
\[
\frac{|f(\gamma(s_n))|}
{\Phi(M_f(|\gamma(s_n)|))}
=
\frac{|f(iy_n)|}{\Phi(M_f(|y_n|))}
\to 0.
\]
Again the desired liminf is zero.

Since \(\gamma\) was arbitrary, Theorem \ref{thm:main} follows.

\section{Consequences and formalization note}

Theorem \ref{thm:main} proves Corollary \ref{cor:third-question} as follows.
Let \(\Psi(T)\to\infty\) be any proposed fixed comparison function.  If
\(\Psi\) is nondecreasing, applying Theorem \ref{thm:main} directly with
\(\Phi=\Psi\) produces a transcendental entire function \(f\) for which every
path to infinity has
\[
\liminf_{t\to\infty}
\frac{|f(\gamma(t))|}{\Psi(M_f(|\gamma(t)|))}=0.
\]
For a nonmonotone \(\Psi\), apply Theorem \ref{thm:main} instead to the
nondecreasing lower envelope \(\Phi\) constructed in the first section.  Since
\(\Phi\le\Psi\) on the relevant tail, any sequence of points for which the
quotient with denominator \(\Phi(M_f(|z|))\) tends to zero also has the quotient
with denominator \(\Psi(M_f(|z|))\) tending to zero.  Hence no path can make the
\(\Psi\)-quotient tend to infinity.  This rules out the universal
comparison-function version of the third question in Erd\H{o}s Problem 514.

The note does not reprove the first two questions in Problem 514, nor does it
claim priority for the parenthetical power-type case discussed in the related
work above.  The direct construction above is self-contained and proves the same
negative statement, while the earlier argument records the shorter derivation
from Hayman's theorem.  The accompanying Lean 4 file
\href{https://github.com/yuta0x89/ErdosProblems/blob/2c0cd8b821ed6ad35c715c0ed67da08bb4494789/Erdos514.lean}{\texttt{Erdos514.lean}}, using
\texttt{leanprover/lean4:v4.28.0} and the corresponding mathlib version,
formalizes the theorem in the equivalent epsilon form.  In that file the main
circle-maximum theorem is \texttt{comparison\_negative\_sphere}; the predicate
\texttt{FrequentlySmall} is the epsilon formulation above; and the circle
maximum \texttt{Ms} matches \(M_f(r)=\max_{|z|=r}|f(z)|\).  The Lean proof first
establishes a closed-disk version, \texttt{comparison\_negative\_closedDisk},
and then identifies it with the circle version by the maximum-modulus principle.
The Lean file treats \(\Phi\) as a total function \(\mathbb R\to\mathbb R\), with
monotonicity and positivity assumed only on \([T_0,\infty)\); this is equivalent
to the partial-domain formulation used here after extending \(\Phi\) arbitrarily
outside the relevant tail.

\section*{Acknowledgements}

The author thanks qrdl for pointing out the relevance of Hayman's theorem and
the surrounding asymptotic-path literature to this question.

\begin{thebibliography}{99}

\bibitem{Erdos1961}
P. Erd\H{o}s,
\emph{Some unsolved problems},
A Magyar Tudom\'{a}nyos Akad\'{e}mia Matematikai Kutat\'{o} Int\'{e}zet\'{e}nek
K\"ozlem\'{e}nyei \textbf{6} (1961), 221--254.

\bibitem{Erdos1982}
P. Erd\H{o}s,
\emph{Some of my favourite problems which recently have been solved},
in \emph{Proceedings of the International Mathematical Conference, Singapore
1981}, North-Holland Math. Stud. \textbf{74}, North-Holland, Amsterdam, 1982,
59--79.

\bibitem{BarthBrannanHayman1978}
K. F. Barth, D. A. Brannan and W. K. Hayman,
\emph{The growth of plane harmonic functions along an asymptotic path},
Proceedings of the London Mathematical Society (3) \textbf{37} (1978),
no. 2, 363--384.

\bibitem{Chojecki2026}
P. Chojecki,
\emph{A note on an Erd\H{o}s path problem for transcendental entire functions},
unpublished note, 2026.

\bibitem{Hayman1960}
W. K. Hayman,
\emph{On the growth of integral functions on asymptotic paths},
The Journal of the Indian Mathematical Society \textbf{24} (1960),
no. 1--2, 251--264.

\bibitem{Huber1957}
A. Huber,
\emph{On subharmonic functions and differential geometry in the large},
Commentarii Mathematici Helvetici \textbf{32} (1957/58), 13--72.

\bibitem{Langley2022}
J. K. Langley,
\emph{Complex flows, escape to infinity and a question of Rubel},
Annales Fennici Mathematici \textbf{47} (2022), no. 2, 885--894.

\bibitem{LewisRossiWeitsman1984}
J. L. Lewis, J. Rossi and A. Weitsman,
\emph{On the growth of subharmonic functions along paths},
Arkiv f\"or Matematik \textbf{22} (1984), 109--119.

\bibitem{RossiWeitsman1985}
J. Rossi and A. Weitsman,
\emph{The growth of entire and harmonic functions along asymptotic paths},
Commentarii Mathematici Helvetici \textbf{60} (1985), 1--14;
correction, ibid. \textbf{60} (1985), 15--16.

\bibitem{Toda1993}
N. Toda,
\emph{A theorem on the growth of entire functions on asymptotic paths and its
application to the oscillation theory of \(w''+Aw=0\)},
Kodai Mathematical Journal \textbf{16} (1993), no. 3, 428--440.

\bibitem{Wu1985}
J.-M. Wu,
\emph{Length of paths for subharmonic functions},
Journal of the London Mathematical Society (2) \textbf{32} (1985),
no. 3, 497--505.

\end{thebibliography}

\end{document}
